% =====================================================================
%  SP519 - Space Data Processing - Rapport de projet
%  Prédiction du flux d'électrons de 1 MeV en orbite basse (NOAA-15)
%  Compilation : pdflatex rapport_SP519.tex   (ou lualatex)
% =====================================================================
\documentclass[11pt,a4paper]{article}

\usepackage[T1]{fontenc}
\usepackage[utf8]{inputenc}
\usepackage[french]{babel}
\usepackage[margin=2cm]{geometry}
\usepackage{amsmath}
\usepackage{booktabs}
\usepackage{graphicx}
\usepackage[hidelinks]{hyperref}

\setlength{\parskip}{0.25em}
\setlength{\parindent}{0pt}

\title{\vspace{-1.6cm}\textbf{Prédiction du flux d'électrons de 1 MeV en orbite basse\\ à partir de NOAA-15 et des données de vent solaire}}
\author{Maxime Cellier \quad\textbullet\quad Alexandre Saint-Paul \\[2pt]
        \small SP519 -- Space Data Processing -- IPSA, automne 2026}
\date{\today}

\begin{document}
\maketitle
\vspace{-0.7cm}

\section{Introduction et objectifs}

Les électrons $>$\,1\,MeV chargent les satellites en orbite basse (LEO) et peuvent endommager
l'électronique ; leur flux varie de plusieurs ordres de grandeur en un à deux jours après une
perturbation du vent solaire. Sa prévision est un enjeu de météorologie de l'espace. NOAA-15 est
un satellite polaire ($\approx 810$\,km) qui traverse les ceintures de radiation à chaque orbite,
la position le long de la ligne de champ étant décrite par le paramètre de McIlwain $L$.

Le projet construit une chaîne complète, des données brutes à la prédiction, sur 2020--2024, avec
deux volets : \textbf{régression} du flux futur à $+12/+24/+36/+48$\,h, et \textbf{classification}
de la survenue d'un « événement » (flux au-dessus du 95\ieme{} percentile).

\section{Données et prétraitement}

\subsection{Sources}
Le flux cible est le canal \texttt{mep\_pro\_tel90\_flux\_p6} de SEM-2 (MEPED) sur NOAA-15
(portail NOAA/NCEI), agrégé heure par heure et par tranche de $L$ de 0,25 : 1\,234\,141 lignes,
sur 1\,808 jours. Ce canal est nominalement un canal \emph{protons} ($>$\,6,9\,MeV) mais, dans
les ceintures, ses comptages sont dominés par les électrons $>$\,1\,MeV — sauf pendant une
éruption solaire à protons, où il mesure à nouveau des protons (cas « où il ne mesure pas des
électrons »).

Les paramètres proviennent d'OMNI2 horaire (NASA SPDF/CDAWeb, lu avec \texttt{cdflib}) : $V$,
$B_z$, $F10.7$, $N$, $P$, $E$, $|B|$ et les indices $K_p$, $Dst$, $AE$ ($<$\,1,6\,\% manquant).

\subsection{Nettoyage et synchronisation}
Sur les \textbf{mesures brutes de 2\,s} : valeurs de remplissage ($-999$) mises à NaN, $L$
invalides et mesures d'étalonnage en vol retirés. Sur la \textbf{table horaire} : une heure sans
mesure ($n=0$) est un trou, un flux nul est au contraire une vraie mesure « 0 particule » ; les
pics aberrants isolés (jusqu'à $9{,}8\times10^{7}$, contre un 99,99\ieme{} percentile
$\approx 10^{6}$) sont retirés au-delà de $10^{7}$ ; les heures contaminées par des protons
solaires sont repérées dans la calotte polaire ($L\ge 10$) par un seuil
\emph{médiane\,+\,10\,$\times$MAD} tenu $\ge 2$\,h, puis retirées. Bilan : 1\,210\,218 lignes
conservées sur 1\,234\,141, dont 23\,923 retirées pour contamination protonique.

\subsection{Région prédite et moyennage sur 12\,h}
La région retenue est $3 \le L < 6$. En dessous de $L=3$, le canal compte surtout des protons
piégés (ceinture interne, anomalie sud-atlantique) ; au-delà de $L\approx 6$, le flux décroît
fortement (574\,\#/cm²/s/sr à $L=6$, 214 à $L=6{,}75$) vers le bruit instrumental
(133\,\#/cm²/s/sr). La région contient le maximum du flux (vers $L\approx 4{,}25$) et remplit les
trois critères attendus : flux $>2\times$ le bruit dans \textbf{92\,\%} des fenêtres, forte
\textbf{variabilité} (écart-type de $\log_{10}$ flux $\approx 0{,}53$), \textbf{lien avec le vent
solaire} (Spearman avec $V$ à $-24$\,h $\approx 0{,}36$).

Tous les paramètres sont moyennés sur des \textbf{fenêtres de 12\,h}. Le flux vaut
$\sum p6\_sum / \sum n$ ; une fenêtre est invalide si elle a moins de 25\,\% du nombre médian de
mesures, et les trous courts ($\le 5$ jours) sont interpolés (colonne \texttt{interpole}). La
table finale compte 3\,654 fenêtres et sert d'entrée à la modélisation.

\section{Analyse statistique}

Le $\log_{10}$ du flux est compact (moyenne $3{,}23$, écart-type $0{,}53$, de $1{,}97$ à
$4{,}64$). Les corrélations \emph{retardées} confirment le délai d'accélération : $V$ atteint
$+0{,}45$ à 60\,h, $K_p$ $+0{,}41$ à 84\,h, $AE$ $+0{,}38$, $Dst$ $-0{,}38$ à 60\,h, $F10.7$
$+0{,}26$ ; $B_z$ instantané est peu corrélé ($-0{,}07$). Les extrêmes des paramètres solaires
tombent sept fois plus souvent pendant un orage (23\,\% contre 3\,\%) : ce sont de vrais
phénomènes, conservés car ce sont les situations à prédire.

Le \textbf{seuil d'événement}, 95\ieme{} percentile du flux des fenêtres mesurées, vaut
\textbf{12\,500\,\#/cm²/s/sr} ($\log_{10}\approx 4{,}10$) : 178 fenêtres (5,0\,\%) en 59 épisodes
(14 en 2024). Pour le volet classification il est \emph{recalculé sur la seule période
d'entraînement}, afin de ne pas utiliser d'information sur le futur.

\section{Modélisation}

\subsection{Protocole}
Le test joue le rôle d'un futur jamais vu : découpage \textbf{chronologique} (entraînement avant
2024, test en 2024), fenêtres interpolées exclues du test. Chaque modèle est comparé à la
\emph{persistance} (flux futur $=$ flux actuel). Les entrées sont le $\log_{10}$ du flux et les
paramètres solaires à $t$, $t-12$\,h et $t-24$\,h.

\subsection{Régression}
Une régression linéaire (entrées standardisées) et une forêt aléatoire sont entraînées par
horizon ; les métriques sur le test (en unités de $\log_{10}$ flux) sont données
tableau~\ref{tab:reg}.

\begin{table}[h]
\centering\small
\caption{Régression : RMSE ($R^2$) sur le test 2024.}
\label{tab:reg}
\begin{tabular}{lcccc}
\toprule
& $+12$\,h & $+24$\,h & $+36$\,h & $+48$\,h \\
\midrule
Persistance          & 0,204 (0,874) & 0,212 (0,863) & 0,289 (0,744) & 0,286 (0,749) \\
Régression linéaire  & \textbf{0,155 (0,927)} & \textbf{0,176 (0,906)} & \textbf{0,216 (0,857)} & \textbf{0,225 (0,844)} \\
Random Forest        & 0,165 (0,917) & 0,196 (0,883) & 0,246 (0,815) & 0,266 (0,782) \\
\bottomrule
\end{tabular}
\end{table}

La régression linéaire \textbf{bat la persistance à tous les horizons}, avec un gain qui croît
avec l'horizon (le vent solaire des jours précédents apporte de l'information que le flux actuel
n'a pas). La forêt fait \emph{moins bien} en moyenne : son erreur se concentre sur les flux bas,
où elle ne peut pas prolonger une tendance ; ailleurs elle est légèrement meilleure. Ses entrées
utiles sont le flux actuel et passé (66 à 83\,\% de l'importance). Sur les \textbf{grandes
variations} (facteur $\ge 3$), les deux modèles réduisent nettement l'erreur de la persistance
(RMSE $\approx 0{,}40$--$0{,}45$ contre $0{,}67$--$0{,}74$), qui reste d'un facteur 2,5 à 3.

\subsection{Classification}
Trois classifieurs sont comparés (régression logistique, SVM RBF, forêt aléatoire ; probabilités
par \texttt{CalibratedClassifierCV} pour le SVM). Le test contient 498 fenêtres dont 40
événements (8,0\,\%) : les classes sont très déséquilibrées et l'\emph{accuracy} seule est
trompeuse (un modèle qui ne prédit jamais d'événement atteindrait déjà 92\,\%). On juge sur l'AUC
et le F1 (tableau~\ref{tab:clf}).

\begin{table}[h]
\centering\small
\caption{Classification : métriques sur le test (498 fenêtres, 40 événements).}
\label{tab:clf}
\begin{tabular}{lccccc}
\toprule
Modèle & Accuracy & Précision & Rappel & F1 & AUC \\
\midrule
Régression logistique & 0,58 & 0,137 & \textbf{0,80} & \textbf{0,234} & \textbf{0,789} \\
SVM (RBF)             & 0,92 & 0,00  & 0,00 & 0,00 & 0,697 \\
Random Forest         & 0,92 & 0,00  & 0,00 & 0,00 & 0,702 \\
\bottomrule
\end{tabular}
\end{table}

Seule la régression logistique détecte réellement des événements. Le SVM et la forêt ont une AUC
proche de 0,70, mais leur seuil implicite de 0,5 sur la probabilité ne produit aucune détection
(classe positive à $\approx 5$\,\%, pondérations de classes). L'AUC, qui mesure la capacité de
\emph{classement}, est donc le résultat le plus robuste.

\section{Discussion et perspectives}

Une part réelle mais modérée du flux est prédictible : la régression linéaire explique 84 à
93\,\% de la variance de $+12$ à $+48$\,h, et l'AUC de classification atteint 0,79. Deux limites
ressortent : la \textbf{taille de l'échantillon d'événements} (le flux de 2024 étant plus bas, le
seuil appris capture davantage de fenêtres au test et le F1 est instable), et la mauvaise
anticipation de la \textbf{frange externe} et des \textbf{orages}.

Pistes : élargir la période de test (2023--2024) et utiliser une validation croisée glissante ;
entraîner des modèles \emph{par tranche de $L$} plutôt qu'une région moyennée ; ajouter des
paramètres et mesurer le gain ; analyser le comportement des modèles pendant la tempête de mai
2024.

\section*{Références}
\small
\begin{itemize}\setlength{\itemsep}{0pt}
  \item NOAA/NCEI, \emph{POES/MetOp SEM-2} : \url{https://www.ncei.noaa.gov/products/poes-metop-space-environment-monitor}
  \item NASA SPDF/CDAWeb, \emph{OMNI2 hourly} : \url{https://spdf.gsfc.nasa.gov/pub/data/omni/omni_cdaweb/hourly/}
  \item Littérature : modèles par tranche de $L$ de type \emph{PreMevE}, ceinture externe
        $2{,}8 \le L \le 6{,}0$ à une résolution de 0,1.
\end{itemize}

\end{document}
